% sec2_8.tex — Section 2.8: Estimation with Parameterized Conditional Heteroskedasticity
%                              (optional)

Even when the error is found to be conditionally heteroskedastic, the OLS estimator is
still consistent and asymptotically normal, and valid statistical inference can be
conducted with robust standard errors and robust Wald statistics.  However, in the
(somewhat unlikely) case of \emph{a priori} knowledge of the functional form of the
conditional second moment $\E(\varepsilon_i^2 \mid \mathbf{x}_i)$, it should be possible
to obtain sharper estimates with smaller asymptotic variance.  Indeed, in finite-sample
theory, the WLS (weighted least squares) can incorporate such knowledge for increased
efficiency in the sense of smaller finite-sample variance.  Does this finite-sample result
carry over to large-sample theory?  This section deals with the large-sample properties
of the WLS estimator.  To simplify the discussion, throughout this section we strengthen
Assumptions~2.2 and 2.5 by assuming that $\{y_i, \mathbf{x}_i\}$ is i.i.d.  This is a
natural assumption to make, because it is usually in cross-section contexts where WLS is
invoked.

%% ─── The Functional Form ─────────────────────────────────────────────────────
\subsection*{The Functional Form}

The parametric functional form for the conditional second moment we consider is
\begin{equation}
  \E(\varepsilon_i^2 \mid \mathbf{x}_i) = \mathbf{z}_i' \bm{\alpha},
  \label{eq:2.8.1}
\end{equation}
where $\mathbf{z}_i$ is a function of $\mathbf{x}_i$.

\begin{example}[parametric form of the conditional second moment]
\label{ex:2.8}
The functional form used in the WLS estimation in the empirical exercise to Chapter~1 was
\[
  \E\!\Bigl(\varepsilon_i^2 \Bigm| \log(Q_i),\, \log(Q_i)^2,\,
    \log\!\bigl(\tfrac{p_{i1}}{p_{i3}}\bigr),\,
    \log\!\bigl(\tfrac{p_{i2}}{p_{i3}}\bigr)\Bigr)
    = \alpha_1 + \alpha_2 \cdot \Bigl(\frac{1}{Q_i}\Bigr).
\]
\end{example}

Since the elements of $\mathbf{z}_i$ can be nonlinear functions of $\mathbf{x}_i$, this
specification is more flexible than it might first look, but still it rules out some
nonlinearities.  For example, the functional form
\begin{equation}
  \E(\varepsilon_i^2 \mid \mathbf{x}_i) = \exp(\mathbf{z}_i' \bm{\alpha})
  \label{eq:2.8.2}
\end{equation}
might be more attractive because its value is guaranteed to be positive.  We consider the
linear specification \eqref{eq:2.8.1} only because estimating parameters in nonlinear
specifications such as \eqref{eq:2.8.2} requires the use of nonlinear least squares.

%% ─── WLS with Known alpha ────────────────────────────────────────────────────
\subsection*{WLS with Known $\bm{\alpha}$}

To isolate the complication arising from the fact that the unknown parameter vector
$\bm{\alpha}$ must be estimated, we first examine the large-sample distribution of the
WLS estimator with known $\bm{\alpha}$.  If $\bm{\alpha}$ is known, the conditional
second moment can be calculated from data as $\mathbf{z}_i' \bm{\alpha}$, and WLS
proceeds exactly as in Section~1.6: dividing both sides of the estimation equation
$y_i = \mathbf{x}_i' \bm{\beta} + \varepsilon_i$ by the square root of
$\mathbf{z}_i' \bm{\alpha}$, to obtain
\begin{equation}
  \tilde{y}_i = \tilde{\mathbf{x}}_i' \bm{\beta} + \tilde{\varepsilon}_i,
  \label{eq:2.8.3}
\end{equation}
where
\[
  \tilde{y}_i \equiv \frac{y_i}{\sqrt{\mathbf{z}_i' \bm{\alpha}}}, \quad
  \tilde{\mathbf{x}}_i \equiv \frac{\mathbf{x}_i}{\sqrt{\mathbf{z}_i' \bm{\alpha}}},
    \quad
  \tilde{\varepsilon}_i \equiv \frac{\varepsilon_i}{\sqrt{\mathbf{z}_i' \bm{\alpha}}},
\]
and then apply OLS.  For later reference, write the resulting WLS estimator as
$\widehat{\bm{\beta}}(\mathbf{V})$.  It is given by
\begin{align}
  \widehat{\bm{\beta}}(\mathbf{V})
    &\equiv \biggl(\frac{1}{n} \sum_{i=1}^{n}
      \tilde{\mathbf{x}}_i \tilde{\mathbf{x}}_i'\biggr)^{\!-1}
      \frac{1}{n} \sum_{i=1}^{n}
      \tilde{\mathbf{x}}_i \cdot \tilde{y}_i \notag \\
    &= \biggl(\frac{1}{n} \sum_{i=1}^{n}
      \frac{1}{\mathbf{z}_i' \bm{\alpha}}\,
      \mathbf{x}_i \mathbf{x}_i'\biggr)^{\!-1}
      \frac{1}{n} \sum_{i=1}^{n}
      \frac{1}{\mathbf{z}_i' \bm{\alpha}}\,
      \mathbf{x}_i \cdot y_i \notag \\
    &= (\mathbf{X}' \mathbf{V}^{-1} \mathbf{X})^{-1}
      \mathbf{X}' \mathbf{V}^{-1} \mathbf{y},
  \label{eq:2.8.4}
\end{align}
with
\[
  \mathbf{V} \equiv \begin{bmatrix}
    \mathbf{z}_1' \bm{\alpha} & & \\
    & \ddots & \\
    & & \mathbf{z}_n' \bm{\alpha}
  \end{bmatrix}.
\]

If Assumption~2.3 is strengthened by the condition that
\begin{equation}
  \E(\varepsilon_i \mid \mathbf{x}_i) = 0,
  \label{eq:2.8.5}
\end{equation}
then $\E(\tilde{\varepsilon}_i \mid \tilde{\mathbf{x}}_i) = 0$.  To see this, note first
that, because $\mathbf{z}_i$ is a function of $\mathbf{x}_i$,
\begin{equation}
  \E(\tilde{\varepsilon}_i \mid \mathbf{x}_i)
    = \E\!\biggl(\frac{\varepsilon_i}{\sqrt{\mathbf{z}_i' \bm{\alpha}}}
      \biggm| \mathbf{x}_i\biggr)
    = \frac{1}{\sqrt{\mathbf{z}_i' \bm{\alpha}}}\,
      \E(\varepsilon_i \mid \mathbf{x}_i) = 0.
  \label{eq:2.8.6}
\end{equation}
Second, because $\tilde{\mathbf{x}}_i$ is a function of $\mathbf{x}_i$, there is no more
information in $\tilde{\mathbf{x}}_i$ than in $\mathbf{x}_i$.  So by the Law of Iterated
Expectations we have
\[
  \E(\tilde{\varepsilon}_i \mid \tilde{\mathbf{x}}_i)
    = \E[\E(\tilde{\varepsilon}_i \mid \mathbf{x}_i) \mid \tilde{\mathbf{x}}_i] = 0.
\]

Therefore, provided that
$\E(\tilde{\mathbf{x}}_i \tilde{\mathbf{x}}_i')$ is nonsingular, Assumptions~2.1--2.5
are satisfied for equation~\eqref{eq:2.8.3}.  Furthermore, by construction, the error
$\tilde{\varepsilon}_i$ is conditionally homoskedastic:
$\E(\tilde{\varepsilon}_i^2 \mid \tilde{\mathbf{x}}_i) = 1$.  So Proposition~2.5 applies
with $\sigma^2 = 1$: the WLS estimator is consistent and asymptotically normal, and the
asymptotic variance is
\begin{align}
  \avar(\widehat{\bm{\beta}}(\mathbf{V}))
    &= \E(\tilde{\mathbf{x}}_i \tilde{\mathbf{x}}_i')^{-1}
    \quad (\text{since the error variance is } 1) \notag \\
    &= \plim\biggl(\frac{1}{n} \sum_{i=1}^{n}
      \tilde{\mathbf{x}}_i \tilde{\mathbf{x}}_i'\biggr)^{\!-1}
    \quad (\text{by ergodic stationarity}) \notag \\
    &= \plim\biggl(\frac{1}{n} \sum_{i=1}^{n}
      \frac{1}{\mathbf{z}_i' \bm{\alpha}}\,
      \mathbf{x}_i \mathbf{x}_i'\biggr)^{\!-1}
    \quad (\text{since } \tilde{\mathbf{x}}_i
      = \mathbf{x}_i / \sqrt{\mathbf{z}_i' \bm{\alpha}}) \notag \\
    &= \plim\biggl(\frac{1}{n} \mathbf{X}' \mathbf{V}^{-1}
      \mathbf{X}\biggr)^{\!-1}.
  \label{eq:2.8.7}
\end{align}
So $\bigl(\frac{1}{n} \mathbf{X}' \mathbf{V}^{-1} \mathbf{X}\bigr)^{-1}$ is a
consistent estimator of $\avar(\widehat{\bm{\beta}}(\mathbf{V}))$.

%% ─── Regression of e_i^2 on z_i ─────────────────────────────────────────────
\subsection*{Regression of $e_i^2$ on $\mathbf{z}_i$ Provides a Consistent Estimate of
$\bm{\alpha}$}

If $\bm{\alpha}$ is unknown, it can be estimated by running a separate regression.
\eqref{eq:2.8.1} says that $\mathbf{z}_i' \bm{\alpha}$ is a regression of
$\varepsilon_i^2$.  If we define
$\eta_i \equiv \varepsilon_i^2 - \E(\varepsilon_i^2 \mid \mathbf{x}_i)$, then
\eqref{eq:2.8.1} can be written as a regression equation:
\begin{equation}
  \varepsilon_i^2 = \mathbf{z}_i' \bm{\alpha} + \eta_i.
  \label{eq:2.8.8}
\end{equation}
By construction, $\E(\eta_i \mid \mathbf{x}_i) = 0$, which, together with the fact that
$\mathbf{z}_i$ is a function of $\mathbf{x}_i$, implies that the regressors
$\mathbf{z}_i$ are orthogonal to the error term $\eta_i$.  Hence, provided that
$\E(\mathbf{z}_i \mathbf{z}_i')$ is nonsingular, Proposition~2.1 is applicable to this
auxiliary regression \eqref{eq:2.8.8}: the OLS estimator of $\bm{\alpha}$ is consistent
and asymptotically normal.

Of course, this we cannot do because we don't observe the error term.  However, since the
OLS estimator $\mathbf{b}$ for the original regression
$y_i = \mathbf{x}_i' \bm{\beta} + \varepsilon_i$ is consistent despite the presence of
conditional heteroskedasticity, the OLS residual $e_i$ provides a consistent estimate of
$\varepsilon_i$.  It is left to you as an analytical exercise to show that, when
$\varepsilon_i$ is replaced by $e_i$ in the regression \eqref{eq:2.8.8}, the OLS
estimator, call it $\hat{\bm{\alpha}}$, is consistent for $\bm{\alpha}$.

%% ─── WLS with Estimated alpha ────────────────────────────────────────────────
\subsection*{WLS with Estimated $\bm{\alpha}$}

The WLS estimation of the original equation $y_i = \mathbf{x}_i' \bm{\beta} +
\varepsilon_i$ with estimated $\bm{\alpha}$ is
$\widehat{\bm{\beta}}(\widehat{\mathbf{V}})$, where $\widehat{\mathbf{V}}$ is the
$n \times n$ diagonal matrix whose $i$-th diagonal is
$\mathbf{z}_i' \hat{\bm{\alpha}}$.  Under suitable additional conditions (see, e.g.,
Amemiya (1977) for an explicit statement of such conditions), it can be shown that
\begin{enumerate}[label=(\alph*)]
  \item $\sqrt{n}\,(\widehat{\bm{\beta}}(\mathbf{V}) - \bm{\beta})$ and
    $\sqrt{n}\,(\widehat{\bm{\beta}}(\widehat{\mathbf{V}}) - \bm{\beta})$ are
    asymptotically equivalent in that their difference converges to zero in probability as
    $n \to \infty$.  Therefore, by Lemma~2.4(a), the asymptotic distribution of
    $\sqrt{n}\,(\widehat{\bm{\beta}}(\widehat{\mathbf{V}}) - \bm{\beta})$ is the same as
    that of $\sqrt{n}\,(\widehat{\bm{\beta}}(\mathbf{V}) - \bm{\beta})$.
    So $\avar(\widehat{\bm{\beta}}(\widehat{\mathbf{V}}))$ equals
    $\avar(\widehat{\bm{\beta}}(\mathbf{V}))$, which in turn is given by
    \eqref{eq:2.8.7};

  \item $\plim \frac{1}{n} \mathbf{X}' \widehat{\mathbf{V}}^{-1} \mathbf{X}
    = \plim \frac{1}{n} \mathbf{X}' \mathbf{V}^{-1} \mathbf{X}$.
    So $\bigl(\frac{1}{n} \mathbf{X}' \widehat{\mathbf{V}}^{-1}
    \mathbf{X}\bigr)^{-1}$ is consistent for
    $\avar(\widehat{\bm{\beta}}(\widehat{\mathbf{V}}))$.
\end{enumerate}

All this may sound complicated, but the operational implication for the WLS estimation of
the equation $y_i = \mathbf{x}_i' \bm{\beta} + \varepsilon_i$ is very clear:

\medskip
\noindent
\textit{Step~1:} \; Estimate the equation $y_i = \mathbf{x}_i' \bm{\beta} +
\varepsilon_i$ by OLS and compute the OLS residuals~$e_i$.

\noindent
\textit{Step~2:} \; Regress $e_i^2$ on $\mathbf{z}_i$, to obtain the OLS coefficient
estimate $\hat{\bm{\alpha}}$.

\noindent
\textit{Step~3:} \; Re-estimate the equation $y_i = \mathbf{x}_i' \bm{\beta} +
\varepsilon_i$ by WLS, using $1/\sqrt{\mathbf{z}_i' \hat{\bm{\alpha}}}$ as the weight
for observation~$i$.

\medskip
\noindent
Since the correct estimate of the asymptotic variance,
$\bigl(\frac{1}{n} \mathbf{X}' \widehat{\mathbf{V}}^{-1} \mathbf{X}\bigr)^{-1}$ in (b),
is ($n$ times) the estimated variance matrix routinely printed out by standard regression
packages for the \textit{Step~3} regression, calculating appropriate test statistics is
quite straightforward: the standard $t$- and $F$-ratios from \textit{Step~3} regression
can be used to do statistical inference.

%% ─── OLS versus WLS ─────────────────────────────────────────────────────────
\subsection*{OLS versus WLS}

Thus we have two consistent and asymptotically normal estimators, the OLS and WLS
estimators.  We say that a consistently and asymptotically normal estimator is
\textbf{asymptotically more efficient} than another consistent and asymptotically normal
estimator of the same parameter if the asymptotic variance of the former is no larger
than that of the latter.  It is left to you as an analytical exercise to show that the WLS
estimator is asymptotically more efficient than the OLS estimator.

The superiority of WLS over OLS, however, rests on the premise that the sample size is
sufficiently large \emph{and} the functional form of the conditional second moment is
correctly specified.  If the functional form is misspecified, the WLS estimator would
still be consistent, but its asymptotic variance may or may not be smaller than
$\avar(\mathbf{b})$.  In finite samples, even if the functional form is correctly
specified, the large-sample approximation will probably work less well for the WLS
estimator than for OLS because of the estimation of extra parameters ($\bm{\alpha}$)
involved in the WLS procedure.
